\textbf{Gene coverage and preprocessing sensitivity.}
All panels use the reference CellPFN checkpoint.
\textbf{a,} Realized genes under the 8,192-HVG cap: means and ranges over three support draws. Labels give panel/supports per class; fewer eligible genes can occur at one-shot support. Source data retain 2k-HVG and all-gene widths.
\textbf{b,} All-gene minus 2k-HVG balanced-accuracy changes, with exported paired means and 95\% $t$ intervals over twenty support draws. Intervals condition on the checkpoint and queries, excluding training and donor-population uncertainty.
\textbf{c,} Random-gene and HVG controls. Small points are twenty-draw means for twelve panel--budget conditions; diamonds equally average conditions. Their spread describes task heterogeneity, not an inferential interval. HVG labels denote caps.
\textbf{d,} Fraction of 21,158 merged cortex genes unmeasured per study and stored as zero. Source data identify support, query and separate sex-case study roles.
\textbf{e,} Changes after removing genes unmeasured in any study, preserving retained processed values and their normalization denominator: paired means and 95\% $t$ intervals over ten draws.
\textbf{f,} Three-draw means for processed cortex versus reconstructed counts. Gene coverage (21,158 versus 12,325) and some support identities differ, preventing an isolated normalization, reconstruction or stored-zero comparison. ``Dev.'' and ``Held'' denote development and held-study cortex; processed cortex remains primary.
